\clearpage\pagestyle{otMain}\setlength{\otBlocEspacement}{2.5mm}\fontsize{11}{15}\selectfont\corrigesChapitre{6}{Additionner et soustraire des fractions}
\begin{blocLivre}{cor-6-1}\begin{corrige}{1}{Addition directe}\demarche{Garder le dénominateur commun.}
On additionne les numérateurs : $(2+3)/7=\res{5/7}$.
\end{corrige}\end{blocLivre}
\begin{blocLivre}{cor-6-2}\begin{corrige}{2}{Soustraction directe}\demarche{Compter les mêmes parts.}
On retranche quatre onzièmes à neuf onzièmes : $\res{5/11}$.
\end{corrige}\end{blocLivre}
\begin{blocLivre}{cor-6-3}\begin{corrige}{3}{Complément à un}\demarche{Écrire l’entier avec le bon dénominateur.}
$1=8/8$, donc $8/8-3/8=\res{5/8}$.
\end{corrige}\end{blocLivre}
\begin{blocLivre}{cor-6-4}\begin{corrige}{4}{Dénominateurs multiples}\demarche{Transformer la première fraction.}
$1/3=2/6$. La somme vaut $3/6=\res{1/2}$.
\end{corrige}\end{blocLivre}
\begin{blocLivre}{cor-6-5}\begin{corrige}{5}{Une part de 28}\demarche{Partager puis prendre les parts.}
Un quart vaut $28:4=7$. Trois quarts valent $7\times3=\res{21}$.
\end{corrige}\end{blocLivre}
\begin{blocLivre}{cor-6-6}\begin{corrige}{6}{Dénominateurs quelconques}\demarche{Choisir un multiple commun.}
Le dénominateur commun 24 donne $15/24+14/24=\res{29/24=1+5/24}$.
\end{corrige}\end{blocLivre}
\begin{blocLivre}{cor-6-7}\begin{corrige}{7}{Une différence}\demarche{Vérifier le sens de la différence.}
Au dénominateur 60 : $42/60-25/60=\res{17/60}$. Le résultat est positif car $7/10>5/12$.
\end{corrige}\end{blocLivre}
\begin{blocLivre}{cor-6-8}\begin{corrige}{8}{Trois termes}\demarche{Utiliser un dénominateur commun à tous.}
Au dénominateur 12 : $8/12+9/12-10/12=\res{7/12}$. Les trois fractions sont réécrites avant d’opérer.
\end{corrige}\end{blocLivre}
\begin{blocLivre}{cor-6-9}\begin{corrige}{9}{Entier et fraction}\demarche{Transformer aussi l’entier.}
Au dénominateur 12 : $36/12+10/12-21/12=\res{25/12=2+1/12}$.
\end{corrige}\end{blocLivre}
\begin{blocLivre}{cor-6-10}\begin{corrige}{10}{Retrouver une part}\demarche{Utiliser l’opération inverse.}
La part manquante vaut $7/6-5/9=21/18-10/18=\res{11/18}$. Contrôle : $10/18+11/18=21/18$.
\end{corrige}\end{blocLivre}
\begin{blocLivre}{cor-6-11}\begin{corrige}{11}{Deux étapes de trajet}\demarche{Additionner les parts du même total.}
On a fait $9/24+10/24=19/24$. Il reste $\res{5/24}$. Une part vaut $144:24=6$ km ; cinq parts valent $\res{30\ \mathrm{km}}$.
\end{corrige}\end{blocLivre}
\begin{blocLivre}{cor-6-12}\begin{corrige}{12}{Un budget}\demarche{Passer de la fraction restante à la somme.}
La part restante vaut $1-8/24-9/24=7/24$. Elle représente $(240:24)\times7=\res{70}$ euros. Il manque $\res5$ euros : l’achat ne peut pas être payé entièrement.
\end{corrige}\end{blocLivre}
\begin{blocLivre}{cor-6-13}\begin{corrige}{13}{Fraction du reste}\demarche{Actualiser la quantité de référence.}
Le premier usage vaut $(90:5)\times2=36$ L. Il reste 54 L. Le second usage vaut $54:3=18$ L. Il reste donc $\res{36\ \mathrm L}$. La seconde fraction porte sur 54 L.
\end{corrige}\end{blocLivre}
\begin{blocLivre}{cor-6-14}\begin{corrige}{14}{Retrouver le total}\demarche{Retrouver la valeur d’une part élémentaire.}
Les deux premières parts valent $3/12+4/12=7/12$. Les 35 cm correspondent aux $5/12$ restants. Un douzième vaut $35:5=7$ cm ; le ruban entier mesure $\res{84\ \mathrm{cm}}$. Contrôle : $21+28+35=84$.
\end{corrige}\end{blocLivre}
\begin{blocLivre}{cor-6-15}\begin{corrige}{15}{Un partage impossible}\demarche{Comparer le total des parts à 1.}
$5/8+7/12=15/24+14/24=29/24>1$. C’est impossible avec un seul gâteau. La promesse dépasse le tout de $\res{5/24}$.
\end{corrige}\end{blocLivre}
